% Values: PRECISE, FAITHFUL, ILLUMINATING, ECONOMICAL
% Evidence planned:
% - PRECISE: a rule for every production of the reduced module grammar
%   (Module, Header in both forms, Imports, ExpList, ExpItem, TopDec in its six
%   forms, DatBind, ConBind, TypBind); every object and operation (ConEnv,
%   Interface, ModEnv, Basis, the merge of environments, the record
%   constructors, the translation of value groups) defined before its first
%   use; one property stated for the interfaces the rules produce.
% - FAITHFUL: imports as elaborated (one qualifier per import, shared
%   qualifiers merge, later imports replace earlier ones, A28, A38); a
%   module's own names replace imported ones; record alias constructors take
%   their fields in the order written; ports bound to outport and inport
%   (D11); each top-level value evaluated once when its module is evaluated,
%   in the dependency order of Appendix A (A59); a delayed self-reference
%   forced during its own evaluation has no result (A25); main need not be
%   exposed (A82); flat modules (D6).
% - ILLUMINATING: comments on the reading of value groups, on why a closed
%   type and an operator item expose nothing, on where a port's payload type
%   comes from, and on the order of imports.
% - ECONOMICAL: options merge the two headers, exposing parts and binding
%   sequences; one rule for a value group, recursive or not; no rule for
%   type-variable sequences or types, which the reduced syntax removes.
\chapter{Dynamic Semantics for Modules}\label{ch:dynamic-modules}

The dynamic semantics of Modules defines the dynamic interfaces that modules
evaluate to, and the environments that imports and top-level declarations
produce. Its sentences infer that a module evaluates to an interface in a
basis, that a sequence of imports evaluates to an environment, that an
exposing list cuts an interface down to the part it exposes, and that a
top-level declaration evaluates to an interface in an environment. Modules
are flat: a module makes its top-level values and constructors available by
name, and there are no signatures or functors. % D6
Every module evaluated here is assumed to have elaborated in the static
semantics of Chapter~\ref{ch:static-modules}. % S4
The order in which the modules of a program are evaluated, and the meaning of
\kw{main}, are given in Chapter~\ref{ch:programs}.

\section{Reduced Syntax}\label{sec:dynamic-modules-reduced}

Types play no part in evaluation, except that a port must know what it sends
or receives. For the dynamic semantics the syntax of Modules is reduced as
follows, in addition to the reduction of the Core
(Section~\ref{sec:dynamic-core-reduced}):
% D1 (definition-style)
\begin{itemize}
\item The type annotations of top-level values are omitted, as in the Core.
  Since a module has no pattern bindings at the top level
  (Section~\ref{sec:syntax-modules-restrictions}), every top-level value
  declaration has the form
  \[
    \var_1\ \kw{=}\ \expm_1\ \kw{and}\ \cdots\ \kw{and}\ \var_n\ \kw{=}\
    \expm_n\qquad(n\geq 1).
  \]
\item The type-variable sequence is omitted from every custom-type binding and
  every type-alias binding; the class TyVarSeq is omitted.
\item Each constructor binding $\con\ \ty_1\cdots\ty_n$ is replaced by $\con$,
  whose arity $\arity(\con)$ is $n$ (Section~\ref{sec:dynamic-core-simple}).
  Thus a constructor binding has the form $\con\ \opt{\kw{|}\ \conbind}$.
\item A type-alias binding $\tycon\ \kw{=}\ \ty$ is replaced by
  $\tycon\ \kw{=}\ \kw{\kwlb}\lab_1\kw{,}\ \cdots\kw{,}\ \lab_n\kw{\kwrb}$
  $(n\geq 0)$ if it is a record alias
  (Section~\ref{sec:syntax-modules-restrictions}) whose record type has the
  labels $\lab_1,\ldots,\lab_n$, in the order written; every other type-alias
  binding is replaced by $\tycon$ alone.
\item A port declaration $\kw{port}\ \var\ \kw{:}\ \ty$ is replaced by
  $\kw{port}\ \var\ \kw{:}\ (\smdirout,\tau)$ for an outgoing port and by
  $\kw{port}\ \var\ \kw{:}\ (\smdirin,\tau)$ for an incoming one, where
  $\tau$ is the payload type of the port; the direction and the payload type
  are those that elaboration records in the port environment (Chapter~\ref{ch:static-modules},
  Section~\ref{sec:static-modules-ports}).
  % DECIDED: the reduced syntax keeps the direction and the payload type
  % tau of a port as elaboration determines them, since encoding and
  % decoding (Chapter 8) are defined by the type, not by the phrase ty.
\item After number specialisation (Section~\ref{sec:dynamic-core-numbers}),
  an exposing item $\var$ naming a top-level variable whose declaration has
  been copied stands for the items naming $\var$ and each of the new
  variables that the copies bind, so that a module that uses $\var$ at an
  instance other than the default one finds the copy for that instance.
  % audit X6, X7: copies other than the default copy bind new variables.
\item An import $\kw{import}\ \modid$ without \kw{as} is read as
  $\kw{import}\ \modid\ \kw{as}\ \modid$: it provides its own \modid{} as
  qualifier (Section~\ref{sec:syntax-modules-identifiers}).
\end{itemize}
As everywhere in the semantics, the derived forms of
Appendix~\ref{ch:appendix-derived-forms} have been applied, so the
top-level declarations of a module stand in the order of
Section~\ref{sec:appendix-derived-forms-modules}; and the imports of a module
are the default imports of Figure~\ref{fig:default-imports} followed by the
imports written in it.

\section{Compound Objects}\label{sec:dynamic-modules-compound}

There are no simple objects beyond those of the Core. The compound objects
for the dynamic semantics of Modules, extra to those of the Core
(Figure~\ref{fig:dynamic-core-compound}), are shown in
Figure~\ref{fig:dynamic-modules-compound}.

\begin{figure}[htbp]
\centering
\begin{objtable}
\objdef{\dmCE}{\dmConEnv}{\cls{TyCon}\finmap\ValEnv}
\objdef{(\dmRA,\dmCE,\VE)\ \text{or}\ I}{\Interface}{\Fin\cls{TyCon}\times\dmConEnv\times\ValEnv}
\objdef{\ME}{\ModEnv}{\cls{ModId}\finmap\Interface}
\objdef{\ME\ \text{or}\ B}{\Basis}{\ModEnv}
\end{objtable}
\caption{Compound Semantic Objects for Modules}
\label{fig:dynamic-modules-compound}
\end{figure}

A \emph{constructor environment} $\dmCE$ maps the type constructor
$\tycon$ of each custom type to a value environment of its constructors. A
dynamic \emph{interface} $(\dmRA,\dmCE,\VE)$ records what a module makes
available to other modules: $\VE$ binds every value it exposes, namely its
exposed values and ports, the constructors of each custom type whose
constructors it exposes, and the constructor of each record alias it
exposes; $\dmCE$ maps every custom type that it exposes to the constructors
of that type that it exposes, which is the empty map for a type exposed
closed; and $\dmRA$ is the set of the type constructors $\tycon$ of the
record aliases whose constructors it exposes. % audit X3, X13
For a type constructor $\tycon$ we write $\con=\tycon$ when $\con$ is spelt
as $\tycon$, as in Chapter~\ref{ch:static-modules}: the constructor of a
record alias $\tycon$ is bound in $\VE$ to the constructor $\con$ with
$\con=\tycon$. A module
environment maps each module identifier to the interface of that module, and
a dynamic basis is a module environment. Interface, ModEnv and Basis are the
dynamic objects of the names that Chapter~\ref{ch:static-modules} uses for
static objects; Chapter~\ref{ch:programs} distinguishes them where both are in
play.

Projection, injection and modification are as for the Core. In particular
``$\dmCE\of I$'' is the second component of $I$; $B(\modid)$ means
$\ME(\modid)$; ``$\VE$ in $\Interface$'' means $(\emptyset,\emptymap,\VE)$;
and interfaces are modified component by component, the set components by
union. The interfaces inferred by the rules of this chapter have the
following property, on which rule~\ruleref{r:dm-expitem-alias} relies.
\begin{property}
If $(\dmRA,\dmCE,\VE)$ is an interface inferred by the rules, then
$\con\in\Dom\VE$ for every $\tycon\in\dmRA$ and $\con=\tycon$, and
$\VE(\con)=v$ whenever $(\dmCE(\tycon))(\con)=v$.
\end{property}

\section{Merging Environments}\label{sec:dynamic-modules-merge}

Several imports of one module may provide the same qualifier, and the names
they provide under it are then all available under that qualifier. Imports
are therefore combined by the following modification, which merges, rather
than replaces, the environment of a qualifier common to both of its
arguments. For environments $E_1=(\QE_1,\VE_1)$ and $E_2=(\QE_2,\VE_2)$,
\[
  E_1\dmmerge E_2 \;=\; (\QE,\ \VE_1+\VE_2),
\]
where $\Dom\QE=\Dom\QE_1\cup\Dom\QE_2$, and $\QE(\modid)$ is
$\QE_1(\modid)+\QE_2(\modid)$ if $\modid\in\Dom\QE_1\cap\Dom\QE_2$, and
otherwise whichever of $\QE_1(\modid)$ and $\QE_2(\modid)$ is defined.
Where both arguments bind a name, the binding of $E_2$ is kept.

\section{Record Constructors}\label{sec:dynamic-modules-record}

The constructor of a record alias is a function of the fields of the record,
taken in the order in which the alias writes them. For $n\geq 1$ and distinct
labels $\lab_1,\ldots,\lab_n$ we extend BasVal with the basic value
$\dmrecord(\lab_1\cdots\lab_n)$, of arity $n$, and extend $\APPLY$ to it by
\[
  \APPLY(\dmrecord(\lab_1\cdots\lab_n),\,v_1\cdots v_n)
  \;=\; \{\{\lab_1\mapsto v_1,\ldots,\lab_n\mapsto v_n\}\}.
\]
For $n=0$ we let $\dmrecord()$ stand for the empty record $\emptymap$, which
is a value but not a basic value, since every basic value has arity at least
1. These values are not bound in the initial basis
(Appendix~\ref{ch:appendix-dynamic-basis}); they are bound only by the
record aliases of modules.

BasVal includes also the values of ports, the basic values
$\effoutport(\var,\tau)$ and $\effinport(\var,\tau)$ for a port name $\var$
and a payload type $\tau$. They are defined in Chapter~\ref{ch:effects}
(Section~\ref{sec:effects-ports}); this chapter uses only the fact that they
are basic values. % D11, audit X11

\section{Top-level Value Groups}\label{sec:dynamic-modules-groups}

A recursive group of top-level values need not consist of functions: the
restrictions of Section~\ref{sec:syntax-modules-restrictions} admit any group
in which every cycle passes through an occurrence within a lambda expression.
% A26
The operation $\Rec$ (Section~\ref{sec:dynamic-core-closures}) gives the
bindings of a group only to the closures that the group binds, not to a
closure held within a record, a tuple or a constructor value. We have
therefore adopted the following translation, under which every member of a
group is reached through a function of the unit value.

Let $\valbind=\var_1\ \kw{=}\ \expm_1\ \kw{and}\ \cdots\ \kw{and}\ \var_n\
\kw{=}\ \expm_n$ be a top-level value declaration, and let
$\var'_1,\ldots,\var'_n$ be distinct variables that occur nowhere in the
program. For a phrase of the Core, $\dmthunk{\mathit{phrase}}$ is the phrase
obtained from it by replacing simultaneously, for $1\leq i\leq n$,
\begin{itemize}
\item each free occurrence of the atomic expression $\var_i$ by the atomic
  expression $\kw{(}\var'_i\ \kw{()}\kw{)}$, and
\item each atomic expression $\kw{\kwlb}\ \var_i\ \kw{|}\ \exprow\
  \kw{\kwrb}$ in which $\var_i$ is free by
  $\kw{let}\ \var_i\ \kw{=}\ \var'_i\ \kw{()}\ \kw{in}\
  \kw{\kwlb}\ \var_i\ \kw{|}\ \dmthunk{\exprow}\ \kw{\kwrb}$,
\end{itemize}
where ``free'' is as in Section~\ref{sec:appendix-derived-forms-let}. These are
the only positions at which a variable occurs free in an expression. Then
\[
  \dmthunk{\valbind}\;=\;
  \var'_1\ \kw{=}\ \kw{\textbackslash}\,\kw{\_}\ \kw{->}\ \dmthunk{\expm_1}\
  \kw{and}\ \cdots\ \kw{and}\
  \var'_n\ \kw{=}\ \kw{\textbackslash}\,\kw{\_}\ \kw{->}\ \dmthunk{\expm_n}.
\]
Every binding of $\dmthunk{\valbind}$ is a lambda expression, so $\Rec$ ties
the knot for it exactly as for a recursive declaration of the Core, and
$\var'_i\ \kw{()}$ then evaluates $\dmthunk{\expm_i}$. When no $\var_i$ is
free in any $\expm_j$, the declaration is not recursive and
$\dmthunk{\expm_i}=\expm_i$.
% DECIDED: a top-level value group is evaluated through this translation.
% Chapter 6's Declarations rules suffice only for groups of lambdas; a
% top-level group may also hold a value that refers to the group within a
% lambda (ones = { head = 1, tail = \_ -> ones }), or a direct reference to
% a member that is a function. The translation keeps every object finite and
% reuses Rec.

\section{Inference Rules}\label{sec:dynamic-modules-rules}

The rules infer sentences of the forms
\[
  \jdg{B}{\mathit{phrase}}{A},\quad
  \jdg{E}{\mathit{phrase}}{A},\quad
  \jdg{I}{\mathit{phrase}}{A},\quad
  \jdg{}{\mathit{phrase}}{A},
\]
where $B$ is a basis, $E$ an environment, $I$ an interface, and $A$ is the
result: a basis with a value environment, a module identifier with an
interface, an environment, an interface, a constructor environment with a
value environment, or a value environment. Each is pronounced ``$\mathit{phrase}$
evaluates to $A$'', followed by ``in $B$'', ``in $E$'' or ``in $I$''. The
hypotheses of a rule that are not sentences are its side conditions; some
premises are sentences of the Core (Chapter~\ref{ch:dynamic-core}). The
convention for options is the same as for the Core static semantics.

As in the Core, sentences have no states and there are no exception packets,
so every rule is given in full, and the order of the premises of a rule
carries no meaning (Section~\ref{sec:dynamic-core-rules}). % D2, A58

\rulegroup{Modules}{\jdg{B}{\modm}{B',\VE}}

\defnrule[r:dm-module]
  {\jdg{B}{\imports}{E}\premsep\jdg{E_0+E}{\topdec}{I}\premsep
   \jdg{I}{\header}{(\modid,I')}}
  {\jdg{B}{\header\ \imports\ \topdec}{\{\modid\mapsto I'\},\ \VE\of I}}

\begin{rulecomments}
\rcomment{r:dm-module} Here $E_0$ is the environment of
Appendix~\ref{ch:appendix-dynamic-basis} that binds the list constructor
names \kw{[]} and \kw{::}; every module starts from it, as its elaboration
starts from $\pi_0$ (rule~\ruleref{r:sm-module}). % audit X4
The top-level declarations are evaluated in
$E_0$ modified by the environment that the imports give, so a name that a module declares replaces,
within it, an imported name spelt alike (rule~\ruleref{r:dm-topdec-seq}). The
basis $\{\modid\mapsto I'\}$ holds only what the header exposes; the value
environment $\VE\of I$ of all the module's own values is the second result,
as in the static semantics, because Chapter~\ref{ch:programs} takes the value
of \kw{main} from it, and \kw{main} need not be exposed. % A82
A module is evaluated once, after the modules it imports, and its top-level
values are evaluated then (Chapter~\ref{ch:programs}). % A59
\end{rulecomments}

\rulegroup{Headers}{\jdg{I}{\header}{(\modid,I')}}

\defnrule[r:dm-header]
  {\jdg{I}{\explist}{I'}}
  {\jdg{I}{\opt{\kw{port}}\ \kw{module}\ \modid\ \kw{exposing}\ \explist}{(\modid,I')}}

\begin{rulecomments}
\rcomment{r:dm-header} The word \kw{port} has no dynamic content: a port
module and a plain module evaluate alike.
\end{rulecomments}

\rulegroup{Imports}{\jdg{B}{\imports}{E}}

\defnrule[r:dm-imports-import]
  {B(\modid)=I\premsep\opt{\jdg{I}{\explist}{I'}}\premsep
   \jdg{B}{\imports}{E}}
  {B\vdash\kw{import}\ \modid\ \kw{as}\ \modid'\
     \opt{\kw{exposing}\ \explist}\ \imports\\
   \Rightarrow(\{\modid'\mapsto(\emptymap,\VE\of I)\},\ \emptymap\opt{+\,\VE\of I'})\dmmerge E}

\defnrule[r:dm-imports-empty]
  {}
  {\jdg{B}{\ }{(\emptymap,\emptymap)}}

\begin{rulecomments}
\rcomment{r:dm-imports-import} $B(\modid)$ is defined, since a module is
evaluated only after every module it imports (Chapter~\ref{ch:programs}). The
qualifier $\modid'$ gives access to every value that the imported module
exposes, its constructors included; the exposing list makes a part of them
available unqualified. The environment of later imports stands on the right of
$\dmmerge$, so a name exposed by a later import replaces the same name exposed
by an earlier one, and in particular a written import replaces a default
import. Where elaboration has found an unqualified name ambiguous, the name is
not used, so which binding is kept is immaterial. % A28, A38
\end{rulecomments}

\rulegroup{Exposing Lists}{\jdg{I}{\explist}{I'}}

\defnrule[r:dm-explist-all]
  {}
  {\jdg{I}{\kw{(..)}}{I}}

\defnrule[r:dm-explist-items]
  {\jdg{I}{\expitem_1}{I_1}\ \ \cdots\ \ \jdg{I}{\expitem_n}{I_n}}
  {\jdg{I}{\kw{(}\expitem_1\ \kw{,}\ \cdots\ \kw{,}\ \expitem_n\kw{)}}
     {I_1+\cdots+I_n}}

\begin{rulecomments}
\rcomment{r:dm-explist-all} The sentence $\jdg{I}{\explist}{I'}$ is the
dynamic counterpart of the interface $\Expose$ of
Chapter~\ref{ch:static-modules}: the values of $I'$ are exactly those whose
names elaboration admits through $\explist$.
\end{rulecomments}

\rulegroup{Exposing Items}{\jdg{I}{\expitem}{I'}}

\defnrule[r:dm-expitem-var]
  {(\VE\of I)(\var)=v}
  {\jdg{I}{\var}{\{\var\mapsto v\}\ \text{in}\ \Interface}}

\defnrule[r:dm-expitem-op]
  {}
  {\jdg{I}{\kw{(}\opv\kw{)}}{(\emptyset,\emptymap,\emptymap)}}

\defnrule[r:dm-expitem-alias]
  {\tycon\in\dmRA\of I\premsep\con=\tycon\premsep(\VE\of I)(\con)=v}
  {\jdg{I}{\tycon}{(\{\tycon\},\emptymap,\{\con\mapsto v\})}}

\defnrule[r:dm-expitem-type]
  {\tycon\notin\dmRA\of I\premsep
   \dmCE=\{\tycon\mapsto\emptymap\,;\,\tycon\in\Dom(\dmCE\of I)\}}
  {\jdg{I}{\tycon}{(\emptyset,\dmCE,\emptymap)}}

\defnrule[r:dm-expitem-open]
  {(\dmCE\of I)(\tycon)=\VE}
  {\jdg{I}{\tycon\,\kw{(..)}}{(\emptyset,\{\tycon\mapsto\VE\},\VE)}}

\begin{rulecomments}
\rcomment{r:dm-expitem-var} The variable is a value or a port.
\rcomment{r:dm-expitem-op} An operator is not in the domain of a value
environment: its value is given by the operator table $\OTdyn$, and is the
same in every module (rule~\ruleref{r:d-atexp-op}). Exposing or importing an
operator therefore binds nothing. % audit X1
\rcomment{r:dm-expitem-alias} Rules~\ruleref{r:dm-expitem-alias}
and~\ruleref{r:dm-expitem-type} cover the one item $\tycon$, distinguished by
whether $\tycon$ is a record alias whose constructor $I$ exposes. The
property of Section~\ref{sec:dynamic-modules-compound} ensures that
$(\VE\of I)(\con)$ is then defined. A record alias exposed by name exposes
its constructor, so \kw{import Geometry exposing (Point)} makes \kw{Point 1 2}
available.
\rcomment{r:dm-expitem-type} A custom type exposed without \kw{(..)}
exposes no value, even when one of its constructors is spelt as the type: in
a module with \kw{type Id = Id Int} and header \kw{exposing (Id)}, the
constructor \kw{Id} is not exposed. The type is still exposed, closed: the
constructor environment maps it to the empty map, so that an import of
\kw{Id(..)} has a meaning (rule~\ruleref{r:dm-expitem-open}). An alias of a
type other than a closed record is in neither $\dmRA\of I$ nor
$\Dom(\dmCE\of I)$, and exposes nothing. % audit X3
\rcomment{r:dm-expitem-open} By elaboration, $\tycon$ names a custom type
that $I$ exposes, so $(\dmCE\of I)(\tycon)$ is defined; an item
$\tycon\,\kw{(..)}$ for an alias is rejected
(Chapters~\ref{ch:syntax-modules} and~\ref{ch:static-modules}). In a header,
$I$ is the module's own interface, which holds every constructor of
$\tycon$. In an import, $\VE$ holds the constructors that the imported
module exposes, and is the empty map when it exposes $\tycon$ closed: with
\kw{module A exposing (T)}, \kw{type T = T Int} and
\kw{import A exposing (T(..))}, the import makes no constructor available.
% audit X3, decision 3
\end{rulecomments}

\rulegroup{Top-level Declarations}{\jdg{E}{\topdec}{I}}

\defnrule[r:dm-topdec-val]
  {\jdg{E}{\dmthunk{\valbind}}{\VE}\\
   \jdg{E+\Rec\VE}{\var'_1\ \kw{()}}{v_1}\ \ \cdots\ \
   \jdg{E+\Rec\VE}{\var'_n\ \kw{()}}{v_n}}
  {\jdg{E}{\valbind}
     {\{\var_1\mapsto v_1,\ldots,\var_n\mapsto v_n\}\ \text{in}\ \Interface}}

\defnrule[r:dm-topdec-type]
  {\jdg{}{\datbind}{(\dmCE,\VE)}}
  {\jdg{E}{\kw{type}\ \datbind}{(\emptyset,\dmCE,\VE)}}

\defnrule[r:dm-topdec-alias]
  {\jdg{}{\typbind}{(\dmRA,\VE)}}
  {\jdg{E}{\kw{type alias}\ \typbind}{(\dmRA,\emptymap,\VE)}}

\defnrule[r:dm-topdec-port-out]
  {}
  {\jdg{E}{\kw{port}\ \var\ \kw{:}\ (\smdirout,\tau)}
     {\{\var\mapsto\effoutport(\var,\tau)\}\ \text{in}\ \Interface}}

\defnrule[r:dm-topdec-port-in]
  {}
  {\jdg{E}{\kw{port}\ \var\ \kw{:}\ (\smdirin,\tau)}
     {\{\var\mapsto\effinport(\var,\tau)\}\ \text{in}\ \Interface}}

\defnrule[r:dm-topdec-seq]
  {\jdg{E}{\topdec_1}{(\dmRA_1,\dmCE_1,\VE_1)}\premsep
   \jdg{E+\VE_1}{\topdec_2}{I_2}}
  {\jdg{E}{\topdec_1\ \kw{;}\ \topdec_2}{(\dmRA_1,\dmCE_1,\VE_1)+I_2}}

\begin{rulecomments}
\rcomment{r:dm-topdec-val} Here $\valbind$ is
$\var_1\ \kw{=}\ \expm_1\ \kw{and}\ \cdots\ \kw{and}\ \var_n\ \kw{=}\ \expm_n$,
and $\dmthunk{\valbind}$ and $\var'_1,\ldots,\var'_n$ are as in
Section~\ref{sec:dynamic-modules-groups}; the first premise is a sentence of
the Core (Value Bindings), and the choice of the new variables does not affect
the result. Each $\var_i$ is bound to a value obtained once, when its module
is evaluated. % A59
A reference to $\var_i$ within the group evaluates $\dmthunk{\expm_i}$ again
in the same environment, which gives $v_i$ again, there being no state,
unless $\expm_i$ relies on the unspecified result of an equality of functions
(Section~\ref{sec:dynamic-core-basic}). A top-level value is thus evaluated
once only up to that non-determinism: where $\expm_i$ compares functions, a
reference to $\var_i$ within its group may give another of the values that
the equality of functions allows, and so differ from $v_i$. % A59 (note), A51, audit X15
Thus with
\kw{ones = \{ head = 1, tail = \textbackslash\_ -> ones \}}, the expression
\kw{(ones.tail ()).head} evaluates to \kw{1}.
% DECIDED: a reference within a group re-evaluates the member; A59 allows
% repeated evaluation where results agree, and only function equality (A51)
% can make them differ.
In \kw{x = let f \_ = x in f ()} the occurrence of \kw{x} is delayed, so the
restrictions of Section~\ref{sec:syntax-modules-restrictions} admit it, but
evaluating $\var'_1\ \kw{()}$ needs $\var'_1\ \kw{()}$ again, without end: no
sentence about the declaration can be inferred, and the module has no
meaning. % A25
\rcomment{r:dm-topdec-port-out} Rules~\ruleref{r:dm-topdec-port-out}
and~\ruleref{r:dm-topdec-port-in}: the payload type $\tau$ comes from
elaboration through the reduced syntax; it has no type variables, so the
value of a port is the same at every use. Applied to a value, the value of an
outgoing port gives a command that sends its encoding at $\tau$; applied to a
function, the value of an incoming port gives a subscription
(Section~\ref{sec:effects-ports}). Port names are unique in a program, so the
name identifies the port. % A84, D11
\rcomment{r:dm-topdec-seq} By the derived forms of
Section~\ref{sec:appendix-derived-forms-modules}, $\topdec$ consists of the
custom-type group, then the aliases, then the ports and value groups in
dependency order, so every declaration is evaluated after each declaration
it refers to. Where that order leaves two declarations unordered, either
order gives the same result.
\end{rulecomments}

\rulegroup{Custom-type Bindings}{\jdg{}{\datbind}{(\dmCE,\VE)}}

\defnrule[r:dm-datbind]
  {\jdg{}{\conbind}{\VE}\premsep\opt{\jdg{}{\datbind}{(\dmCE,\VE')}}}
  {\jdg{}{\tycon\ \kw{=}\ \conbind\ \opt{\kw{and}\ \datbind}}
     {(\{\tycon\mapsto\VE\}\opt{+\,\dmCE},\ \VE\opt{+\,\VE'})}}

\begin{rulecomments}
\rcomment{r:dm-datbind} No type names are generated: a custom type has no
dynamic content beyond its constructors. The constructors of a module are
distinct (Section~\ref{sec:syntax-modules-restrictions}), so the
modification of $\VE$ replaces no binding.
\end{rulecomments}

\rulegroup{Constructor Bindings}{\jdg{}{\conbind}{\VE}}

\defnrule[r:dm-conbind]
  {\opt{\jdg{}{\conbind}{\VE}}}
  {\jdg{}{\con\ \opt{\kw{|}\ \conbind}}{\{\con\mapsto\con\}\opt{+\,\VE}}}

\begin{rulecomments}
\rcomment{r:dm-conbind} The constructor value $\con$ awaits all
$\arity(\con)$ of its arguments; a constructor of arity 0 is thus a complete
value. Constructors spelt alike in different modules give equal constructor
values, which never meet in one rule, since they have different types
(Section~\ref{sec:dynamic-core-simple}).
\end{rulecomments}

\rulegroup{Alias Bindings}{\jdg{}{\typbind}{(\dmRA,\VE)}}

\defnrule[r:dm-typbind]
  {\opt{\con=\tycon}}
  {\vdash\tycon\ \opt{\kw{=}\ \kw{\kwlb}\lab_1\kw{,}\ \cdots\kw{,}\ \lab_n\kw{\kwrb}}\\
     \Rightarrow(\emptyset\opt{\cup\,\{\tycon\}},\ \emptymap\opt{+\,\{\con\mapsto\dmrecord(\lab_1\cdots\lab_n)\}})}

\begin{rulecomments}
\rcomment{r:dm-typbind} Here $n\geq 0$, and by the syntactic restrictions
the labels are distinct. The constructor takes the fields in the order in
which the alias writes them: with \kw{type alias P = \{ y : Int, x : Int \}},
\kw{P 1 2} evaluates to \kw{\{ y = 1, x = 2 \}}. An alias of the empty record
binds its name to the empty record itself
(Section~\ref{sec:dynamic-modules-record}). An alias of any other type binds
nothing, and its type constructor is not in $\dmRA$; so rule
\ruleref{r:dm-topdec-alias} records in $\dmRA$ exactly the record aliases.
% audit X13
\end{rulecomments}
